% !TeX root = ../main.tex

\chapter{Introduzione e Definizione del Servizio}
\label{cap:introduzione_definizione}

% ---------------------------------------------------------
\section{Dominio Operativo}
\label{sec:dominio_operativo}

Il progetto si colloca all'intersezione tra due domini critici identificati dalle specifiche di progetto: \textit{Transportation} e \textit{Industrial / IoT Monitoring}. % 
Nello specifico, il sistema implementa un'architettura logistica autonoma per la gestione di una flotta di droni adibita alla consegna urbana di \textit{payload} critici (come forniture mediche d'emergenza, organi per trapianti o componentistica industriale ad alta priorità). % 
In un servizio "high-stakes" come questo l'efficienza e l'affidabilità del sistema sono fondamentali per la buona riuscita di missioni. % 

\subsection*{Il Contesto Fisico e Simulato}

Il sistema opera all'interno di un'area metropolitana modellata metricamente, con un raggio d'azione di 5000 metri (5 km) attorno a un HUB logistico centrale posizionato alle coordinate base \texttt{[0.0, 0.0]}. % 
L'ambiente operativo è altamente dinamico e soggetto a vincoli fisici realistici che l'Intelligenza Artificiale deve interpretare e gestire: % 

\begin{itemize}
    \item \textbf{Gestione del Payload:} Gli ordini generati presentano un peso variabile (da 0.5 kg a 5.0 kg) e tre livelli di priorità (\textit{low, normal, high}). % 
    \item \textbf{Fisica del Volo:} I droni non sono entità ideali, ma agenti fisici il cui stato viene modellato in base allo sforzo. % 
    Il consumo della batteria e la velocità di crociera in consegna (che varia dai 50 km/h agli 80 km/h) sono calcolati dinamicamente in proporzione diretta al peso del carico trasportato; in fase di rientro senza carico il simulatore può arrivare fino a 100 km/h. % 
    Il drone può trovarsi in stato IDLE, IN DELIVERY, RETURNING e MAINTENANCE. % 
    \item \textbf{Degrado Meccanico (Wear):} Ogni missione accumula una percentuale di usura (da 0\% a 100\%) basata sulla distanza percorsa. % 
    Il superamento di soglie critiche richiede interventi di manutenzione preventiva e ragionamenti per evitare schianti catastrofici. % 
\end{itemize}

\subsection*{La Necessità di un Approccio Agentico}

In un dominio tradizionale, l'assegnazione degli ordini avverrebbe tramite semplici code FIFO (First-In, First-Out) o algoritmi deterministici di routing. % 
Tuttavia, l'elevata variabilità delle condizioni (es. un pacco da 4.5 kg ad altissima priorità, ma con a disposizione solo droni con batteria bassa o con usura avanzata) rende l'automazione rigida inefficiente e pericolosa. % 
Il sistema delega questo triage a un layer operativo basato su Agentic AI. % 
L'obiettivo del dominio applicativo è dimostrare come l'IA possa orchestrare un'infrastruttura IoT in tempo reale, bilanciando il rischio operativo, l'urgenza delle consegne e la salute della flotta. % 

% ---------------------------------------------------------
\section{Stakeholder e Utenti Target}
\label{sec:stakeholder_utenti}

Per comprendere appieno i requisiti di affidabilità, sicurezza e scalabilità dell'architettura, è necessario definire gli attori che interagiscono con i vari componenti. % 
Gli attori si dividono in due macro-categorie: gli \textit{stakeholder} (coloro che hanno un interesse diretto nella gestione e sostenibilità economica della piattaforma) e gli \textit{utenti target} (i beneficiari del servizio logistico). % 

\subsection*{Stakeholder}

\begin{enumerate}
    \item \textbf{Operatori di Flotta e Supervisori Umani (Human-in-the-Loop):} 
    Sono gli utenti responsabili della supervisione strategica e della sicurezza operativa del sistema. % 
    Interagiscono direttamente con il gateway di sicurezza per validare o respingere le azioni ad alto impatto generate autonomamente dall'Intelligenza Agentica. % 
    Tra queste azioni figurano lo scaling della flotta oltre le soglie di sicurezza automatizzate o l'invio di comandi MQTT critici. % 
    Gli operatori monitorano inoltre lo stato in tempo reale dei droni e degli ordini tramite la dashboard web. % 

    \item \textbf{Ingegneri di Infrastruttura e DevOps/SRE Team:}
    Professionisti incaricati di garantire l'alta disponibilità e l'affidabilità dell'architettura. % 
    Il loro obiettivo principale è monitorare la salute dei nodi e dei pod Kubernetes che ospitano i droni simulati. % 
    Verificano inoltre il corretto instradamento dei messaggi sul broker Mosquitto e assicurano la continuità della persistenza dei dati su InfluxDB. % 
    Intervengono in caso di escalation infrastrutturale qualora gli agenti riproducano log di errore legati a guasti fisici o problemi di risoluzione DNS non mitigabili autonomamente. % 

    \item \textbf{Business Owner e Controller Finanziari:}
    Soggetti interessati alla sostenibilità economica e all'efficienza della piattaforma. % 
    Valutano il bilanciamento tra investimenti infrastrutturali (CAPEX) e costi operativi (OPEX). % 
    Pongono particolare attenzione al monitoraggio del tetto massimo di spesa legato al consumo di token per l'inferenza dei modelli LLM (\textit{Agent-related cost ceiling}). % 
    Il loro KPI di riferimento è il \textit{Return on Agent} (ROA). % 
    Questo indicatore misura il risparmio derivante dalla prevenzione dei guasti hardware tramite manutenzione predittiva rispetto al costo vivo delle chiamate API dell'IA. % 
\end{enumerate}

\subsection*{Utenti Target (Clienti del Servizio)}

\begin{itemize}
    \item \textbf{Entità Mittenti (Hub Medici, E-Commerce Critici, Hub Industriali):}
    Rappresentano i clienti primari del software logistico. % 
    Sono responsabili dell'immissione asincrona di richieste di consegna all'interno della rete. % 
    Questi utenti richiedono il rigido rispetto dei \textit{Service Level Objectives} (SLO) concordati. % 
    Esigono che gli ordini ad alta priorità (\textit{high priority}) vengano presi in carico il prima possibile e consegnati con tempistiche accettabili anche in scenari di forte congestione della coda (\textit{load spikes}). % 

    \item \textbf{Destinatari Finali (Laboratori di Analisi, Linee di Produzione JIT):}
    I beneficiari ultimi del payload trasportato dai velivoli autonomi. % 
    Nel contesto di spedizioni urgenti o industriali (\textit{Just-In-Time}), il destinatario dipende strettamente dalla precisione e dalla resilienza del sistema. % 
    Per questa categoria di utenti, la minimizzazione del downtime e l'eliminazione dei fallimenti (es. schianto del drone per esaurimento della batteria in volo) rappresentano requisiti di affidabilità critici. % 
\end{itemize}

% ---------------------------------------------------------
\section{Analisi del Workload}
\label{sec:analisi_workload}

L'architettura è di tipo \textit{event-driven}, e sottoposta a un carico di lavoro ibrido, caratterizzato dalla coesistenza di flussi di dati continui (streaming IoT) e da eventi asincroni variabili (richieste di business e cicli di ragionamento AI). Per dimensionare correttamente i meccanismi di reliability e scalabilità, il workload è stato scomposto in tre sezioni principali:

\subsection*{1. Telemetria IoT ad Alta Frequenza (Constant Load)}
Il carico di base dell'infrastruttura è generato dalla flotta. Ogni nodo (pod) del \texttt{drone\_simulator} calcola costantemente il proprio movimento e il decadimento della batteria, trasmettendo un payload JSON contenente coordinate spaziali, stato operativo e usura. 
Ciascun drone pubblica un aggiornamento sul broker MQTT (topic \texttt{telemetry/drones}) con una frequenza di un messaggio ogni 2 secondi. Questo genera un throughput in ingresso costante, che il \texttt{central\_server} deve prima elaborare in RAM per mantenere aggiornata la Dashboard real-time, e successivamente persistere su InfluxDB strutturando i dati come serie temporali.

\subsection*{2. Eventi di Business asincroni (Spiky Load)}
Il secondo vettore è rappresentato dalle richieste di consegna generate dai nodi cliente (\texttt{client\_simulator.py}). Questo traffico è composto dagli ordini generati casualmente e emessi sul topic \texttt{business/ordini/nuovi}. Allo stesso tempo il sistema è soggetto a picchi anomali in scenari di stress testing che forzano l'intervento reattivo del sistema. 

\subsection*{3. Carico di Inferenza e Polling Agentico (Read-Heavy Spikes)}
Una peculiarità delle architetture che sfruttano AI Agentica è il carico indiretto generato dall'Intelligenza Artificiale stessa sull'infrastruttura sottostante. L'orchestratore centrale (\texttt{logistic\_ai\_brain.py}) si occupa di risvegliare gli agenti e ogni risveglio genera varie operazioni in lettura (Polling):
\begin{itemize}
    \item Interrogazione intensiva di InfluxDB per estrarre la finestra temporale (time-window) degli ultimi 5-60 minuti di telemetria e ordini.
    \item Chiamate alle API di Kubernetes per ottenere lo stato dei Pod attivi in tempo reale.
\end{itemize}
Questo workload "a raffiche" richiede che il database e le API di K8s rispondano con latenze minime, poiché un ritardo nella fase di \textit{information gathering} prolungherebbe i tempi di inferenza del Large Language Model, violando i \textit{Service Level Objectives} (SLO) di reattività del sistema. % 

% ---------------------------------------------------------